\noindent\printedpage{54}\msp{page 108}
% head: A. Grothendieck : Dérivateurs, Ch. XVII, Édition des Universités de Montpellier II et Paris VII
\[
\begin{tikzcd}
X(1) \arrow[r, "X(\partial)"] \arrow[d, "u(1)"'] \arrow[dr, phantom, "\text{cocart.}"] & X(0) \arrow[d] \arrow[dd, dotted, bend left=40, "u(0)"] \\
Y(1) \arrow[r] \arrow[dr, dotted, "Y(\partial)"'] & \tilde{Y} \arrow[d, "\tilde{u}"'] \\
 & Y(0)
\end{tikzcd}
\qquad .
\tag*{(6.2.5.4)}
\]

For us to have $u \in \Phi(\boldsymbol{\Delta}^1)$, it is necessary and sufficient that we have
\[
u(0),\ u(1),\ \tilde{u} \ \in\ \Phi\ .
\tag*{(6.2.5.5)}
\]

I have already checked this (in dual form), but can no longer find the place (perhaps I did not take notes). I don't feel like looking for the proof here.

To finish proving 6.1. d), i.e. the equality $C(I) = C_I$ for $I$ artinian, I feel like proving something still more general, as follows: % [sic] plus générale -> plus général, agreement with « quelque chose »

\textbf{Corollary 6.3.} \textit{Let $(\mathcal{M}, \Phi, \Psi)$ and $I \longmapsto \Phi(I) \subset \mathrm{Fl}(\mathcal{M}(I))$ be as in \textup{6.2.5.} Let $I$ be an ordered set, and let us also consider $\Phi_I$ defined (in the dual situation) on page \textup{25 (2.2)} (but with dual notation \dots). I say that we have}
\[
\Phi_I \ \supset\ \Phi(I)\ ,
\tag*{(6.3.1)}
\]
\msp{page 109}\textit{and we have equality if $I$ is noetherian, i.e. $I^{\circ}$ is artinian.}

\textsc{Proof}. The proof of $\Phi(I) \subset \Phi_I$ is the one we have just given. It remains to prove that if $I$ is noetherian, we have $\Phi_I \subset \Phi(I)$. Given $i : A \longrightarrow B$ in $\Phi_I$, we must prove that $u \in \Phi(I)$, % [sic] u -> i, the arrow in question is i : A -> B
i.e. that for a commutative diagram in $\mathcal{M}(I)$,
\[
\begin{tikzcd}
A \arrow[r, "\alpha"] \arrow[d, "i"'] & X \arrow[d, "p"] \\
B \arrow[r, "\beta"'] \arrow[ur, dotted, "h"] & Y
\end{tikzcd}
\tag*{(6.3.2)}
\]
with $p \in \Psi\langle I\rangle$, i.e. $p(i) : X(i) \longrightarrow Y(i)$ in $\Psi$ for every $i$, $\exists\ h : B \longrightarrow X$ leaving the diagram commutative. We are going to return to the dual notation, working in $\mathcal{M}(I^{\circ})$, with $I$ artinian. We consider the set $\mathcal{H}$ of pairs % [sic] des couple -> des couples, plural
$(U, h_U)$ of an open set $U$ of $I$ and an $h_U : B|U \longrightarrow X|U$, making the restriction to $U$ of (6.3.2) commutative. % [sic] commutatif -> commutative, agreement with « la restriction »
We order $\mathcal{H}$ by extension, we note that every totally ordered subset $(U_\alpha, h_{U_\alpha})$ in $\mathcal{H}$ has a least upper bound, namely $U = \bigcup_\alpha U_\alpha$, $h_U$ defined by the $h_{U_\alpha}$ on the $U_\alpha$.
% foot: Version 9 février 2019

\printedpage{55}\msp{page 110}
% head: A. Grothendieck : Dérivateurs, Ch. XVII, Édition des Universités de Montpellier II et Paris VII
(\textbf{N.B.} The set $\mathcal{H}$ is not empty, for $(\emptyset, h_{\emptyset})$ is in $\mathcal{H}$ -- but this is contained in the preceding statement that every totally ordered subset of $\mathcal{H}$ has a least upper bound in it \dots). By Zorn, there exists a maximal element $(U, h_U)$. To prove that $U = I$, and we will be done. Or again that if $(U, h_U)$ such that $U \neq I$, it is \textit{not} maximal. Let $F = I\backslash U \neq \emptyset$ be the complementary closed set. By the artinian hypothesis on $I$, $F$ has a minimal element $i_0$. Let
\[
U' \ =\ U \cup \overbrace{I/i_0}^{V} \ =\ U \cup \{i_0\}\ ,
\]
we are going to extend $h_U$ to $h_{U'}$ on $U'$, we will be done. For this, let $V_0 = V \cap U = V\backslash\{i_0\}$, we are reduced to extending the restriction $h_{V_0}$ of $h_U$ to $V_0$, from $V_0$ to $V$. All told, we may suppose $I = I/i_0$ equipped with a final element $i_0$, $U = I\backslash\{i_0\}$, and show that a $(U, h_U) \in \mathcal{H}$ extends to $I$. Let $A'$, $B'$, $X'$, $Y'$ be the restrictions of $A$, $B$, $X$, $Y$ to $U$ (which we may also denote by $I'$). The datum of $A$ % [sic] La donné -> La donnée, missing feminine ending
is equivalent to that of $A'$, that of $A_0 = A(i_0) \in \mathrm{Ob}\,\mathcal{M}$, and of an arrow from
\[
A'_{I'} \ \overset{\text{def}}{=}\ \varinjlim\limits_{i' \in I'} A'(i') \ \xrightarrow{\ \lambda_A\ }\ A_0\ (\ \simeq A_I\ )\ ,
\]
\msp{page 111}and likewise for $B$, $X$, $Y$. More precisely, the category $\mathcal{M}(I^{\circ})$ can be interpreted as that of triples $X = (X', X_0, \lambda_X)$, where $X' \in \mathcal{M}(I'^{\circ})$, $X_0 \in \mathcal{M}$, and $\lambda_X : X'_{I'} \longrightarrow X_0$. In these terms, the datum of the diagram (6.3.2) is equivalent to that of a similar square diagram
\[
\begin{tikzcd}
A' \arrow[r, "\alpha'"] \arrow[d, "i'"'] & X' \arrow[d, "p'"] \\
B' \arrow[r, "\beta'"'] \arrow[ur, dotted, "h'"] & Y'
\end{tikzcd}
\tag*{(6.3.3)}
\]
\textit{in} $\mathcal{M}(I'^{\circ})$, whence a diagram in $\mathcal{M}$
\[
\begin{tikzcd}
A'_{I'} \arrow[r, "\alpha'_{I'}"] \arrow[d, "i'_{I'}"'] & X'_{I'} \arrow[d, "p'_{I'}"] \\
B'_{I'} \arrow[r, "\beta'_{I'}"'] \arrow[ur, dotted, "h'_{I'}"] & Y'_{I'}\ ,
\end{tikzcd}
\tag*{(6.3.4)}
\]
of another square diagram \textit{in} $\mathcal{M}$
\[
\begin{tikzcd}
A_0 \arrow[r, "\alpha_0"] \arrow[d, "i_0"'] & X_0 \arrow[d, "p_0"] \\
B_0 \arrow[r, "\beta_0"'] & Y_0\ ,
\end{tikzcd}
\tag*{(6.3.5)}
\]
and finally of a homomorphism $\lambda = (\lambda_A, \lambda_B, \lambda_X, \lambda_Y)$ from the square (6.3.4) to the square (6.3.5):
% foot: Version 9 février 2019
\printedpage{56}\msp{page 112}
% head: A. Grothendieck : Dérivateurs, Ch. XVII, Édition des Universités de Montpellier II et Paris VII
\[
\begin{tikzcd}
A'_{I'} \arrow[rr, "\alpha'_{I'}"] \arrow[dr, "i'_{I'}"'] \arrow[dd, "\lambda_A"'] & & X'_{I'} \arrow[dr, "p'_{I'}"] \arrow[dd, "\lambda_X" near start] & \\
 & B'_{I'} \arrow[ur, dotted, "h'_{I'}"] \arrow[rr, crossing over, "\beta'_{I'}" near end] & & Y'_{I'} \arrow[dd, "\lambda_Y"] \\
A_0 \arrow[rr, "\alpha_0" near start] \arrow[dr, "i_0"'] & & X_0 \arrow[dr, "p_0"] & \\
 & B_0 \arrow[from=2-2, crossing over, "\lambda_B"' near start] \arrow[ur, dotted, "h_0"'] \arrow[rr, "\beta_0"] & & Y_0
\end{tikzcd}
\qquad\qquad \text{\footnotemark[32]}
\tag*{(6.3.6)}
\]
\footnotetext[32]{\textbf{N.B.} The condition $p \in \Psi\langle I\rangle$ is expressed by $p' \in \Psi\langle I'\rangle$ and \fbox{$p_0 \in \Psi.$}}

On the other hand, the question of extending $(I', h') \in \mathcal{H}$ to $(I, h)$, amounts to that of finding $h_0 : B_0 \longrightarrow X_0$ in $\mathcal{M}$, so that together with $h'_{I'}$ this again makes a commutative diagram of (6.3.6), which amounts to saying that on the one hand the two triangles formed in the base square are commutative:
\[
p_0h_0 \ =\ \beta_0, \qquad\quad h_0i_0 \ =\ \alpha_0\ ,
\tag*{(6.3.7)}
\]
and on the other hand the diagonal square defined by $h'_{I'}$, $h_0$ is commutative (which expresses that we have an arrow $h$ from $B$ to $X$):
\[
\lambda_X h'_{X'} \ =\ h_0\lambda_B\ . % [sic] h'_{X'} -> h'_{I'}, the arrow is h'_{I'} everywhere else (cf. (6.3.16))
\tag*{(6.3.8)}
\]

\msp{page 113}But the diagram (6.3.6), where the densely dotted arrow $h'_{I'}$ is already in hand, and $h_0$ in spaced dots remains to be found, gives rise to the diagram, in which is brought out the arrow $\tilde{i} \in \Phi$ which expresses the hypothesis $i \in C(i)$, % [sic] C(i) -> C(I), cf. i' \in C(I')
in view of $i' \in C(I')$:
\[
\begin{tikzcd}
A'_{I'} \arrow[rr, "\alpha'_{I'}"] \arrow[dr, "i'_{I'}"'] \arrow[dd, "\lambda_A"'] & & X'_{I'} \arrow[dr, "p'_{I'}"] \arrow[dd, "\lambda_X" near start] & \\
 & B'_{I'} \arrow[ur, dotted, "h'_{I'}"] \arrow[rr, crossing over, "\beta'_{I'}" near end] & & Y'_{I'} \arrow[dd, "\lambda_Y"] \\
A_0 \arrow[ur, phantom, "\text{cocart.}"] \arrow[rr, "\alpha_0" near start] \arrow[dr, "r"'] \arrow[ddr, dotted, "i_0"'] & & X_0 \arrow[dr, "p_0\in\Psi"] & \\
 & \tilde{B} \arrow[from=2-2, crossing over, "s"' near start] \arrow[ur, dotted, "\tilde{h}_0"] \arrow[rr, dotted, "\tilde{\beta}_0"] \arrow[d, "\tilde{i}\in\Phi"'] & & Y_0 \\
 & B_0\ , \arrow[uur, dotted, "h_0"'] \arrow[urr, "\beta_0"'] & &
\end{tikzcd}
\tag*{(6.3.9)}
\]
% foot: Version 9 février 2019
\printedpage{57}
% head: A. Grothendieck : Dérivateurs, Ch. XVII, Édition des Universités de Montpellier II et Paris VII
where the densely dotted arrow $\tilde{h}_0$ from $\tilde{B} \overset{\text{def}}{=} A_0 \amalg_{A'_{I'}} B'_{I'}$ to $X_0$, is defined by the commutativity conditions
\[
\tilde{h}_0r \ =\ \alpha_0, \qquad\quad \tilde{h}_0s \ =\ \lambda_X h'_{I'}\ ,
\tag*{(6.3.10)}
\]
which is legitimate, seeing that
\[
\begin{aligned}
\alpha_0\lambda_A \ &=\ (\lambda_X h'_{I'})i'_{I'} \\
&=\ \lambda_X(h'_{I'}i'_{I'}) \\
&=\ \lambda_X\alpha'_{I'}.
\end{aligned}
\]

\msp{page 114}As for the arrow $\tilde{\beta}_0 : \tilde{B} \longrightarrow Y_0$, it is given by the commutativity condition
\[
\tilde{\beta}_0r \ =\ p_0\alpha_0, \qquad\quad \tilde{\beta}_0s \ =\ \lambda_Y\beta'_{I'}\ ,
\tag*{(6.3.11)}
\]
which is again legitimate, since
\[
(p_0\alpha_0)\lambda_A \ =\ (\lambda_Y\beta'_{I'})i'_{I'}
\]
by virtue of the commutativity of the cube (6.3.6). To establish the full commutativity of the diagram (6.3.9) (solid lines, and densely dotted lines for $h'_{I'}$, $\tilde{h}_0$ and $\tilde{\beta}_0$ -- excluding $h_0$ not yet constructed), we must still verify on the one hand the relations % [sic] les relation -> les relations, plural
\[
\left\{
\begin{aligned}
\tilde{h}_0r \ &=\ \alpha_0, \quad p_0\tilde{h}_0 \ =\ \tilde{\beta}_0, \\
\lambda_X h'_{I'} \ &=\ \tilde{h}_0s
\end{aligned}
\right.
\qquad\qquad \text{\footnotemark[33]},
\tag*{(6.3.12)}
\]
\footnotetext[33]{\textbf{N.B.} the first and the third of these relations are already known, cf. (6.3.10).}
establishing the commutativity of the upper cube with the diagonal square (diagram of type (6.3.6), the dotted arrow $h_0$ included in it), on the other hand the relations
\[
\left\{
\begin{aligned}
\tilde{i}r \ &=\ i_0 \qquad \text{(and $\tilde{i}s = \lambda_B$ for the record)} \\
\beta_0\tilde{i} \ &=\ \tilde{\beta}_0\ ,
\end{aligned}
\right.
\tag*{(6.3.13)}
\]
where the first line is already known by the definition of $\tilde{B}$ and of $\tilde{i}$. The second
\msp{page 115}reduces to the two relations
\[
\begin{aligned}
(\beta_0\tilde{i})r \ &=\ \beta_0(\tilde{i}r) \\
&=\ \beta_0i_0 \\
&=\ p_0\alpha_0 \\
&=\ \tilde{\beta}_0r\ ,
\end{aligned}
\]
OK by commutativity of the base square of (6.3.5),
\[
\begin{aligned}
(\beta_0\tilde{i})s \ &=\ \beta_0(\tilde{i}s) \\
&=\ \beta_0\lambda_B \\
&=\ \lambda_Y\beta'_{I'} \\
&=\ \tilde{\beta}_0s\ ,
\end{aligned}
\]
% foot: Version 9 février 2019
\printedpage{58}
% head: A. Grothendieck : Dérivateurs, Ch. XVII, Édition des Universités de Montpellier II et Paris VII
OK by commutativity of the front square in (6.3.6).

So it remains to verify the second relation (6.3.12),
\[
p_0\tilde{h}_0 \ =\ \tilde{\beta}_0\ ,
\]
which reduces to
\[
\begin{aligned}
(p_0\tilde{h}_0)r \ &=\ p_0(\tilde{h}_0r) \\
&=\ p_0\alpha_0 \\
&=\ \tilde{\beta}_0r\ ,
\end{aligned}
\]
OK by (6.3.11),
\[
\begin{aligned}
(p_0\tilde{h}_0)s \ &=\ p_0(\tilde{h}_0s) \\
&=\ p_0(\lambda_X h'_{I'}) \\
&=\ (p_0\lambda_X)h'_{I'} \\
&=\ (\lambda_Y p'_{I'})h'_{I'} \\
&=\ \lambda_Y(p'_{I'}h'_{I'}) \\
&=\ \lambda_Y\beta'_{I'} \\
&=\ \tilde{\beta}_0s\ ,
\end{aligned}
\]
OK.

This being settled, I say that the constraints (6.3.7), (6.3.8) on $h_0$ amount to the commutativity of the two triangles in the square % [sic] les carré -> le carré, singular (one square (6.3.14))
\[
\begin{tikzcd}
\tilde{B} \arrow[r, "\tilde{h}_0"] \arrow[d, "\tilde{i}\in\Phi"'] & X_0 \arrow[d, "p_0\in\Psi"] \\
B_0 \arrow[r, "\beta_0"'] \arrow[ur, dotted, "h_0"] & Y_0\ ,
\end{tikzcd}
\tag*{(6.3.14)}
\]
\msp{page 116}i.e. to the two relations
\[
h_0\tilde{i} \ =\ \tilde{h}_0\ , \qquad\quad p_0h_0 \ =\ \beta_0\ .
\tag*{(6.3.15)}
\]

The second of these relations is none other than the first relation (6.3.7), and it remains to see that the first relation (6.3.15) is equivalent to the two remaining relations in (6.3.7), (6.3.8) taken together:
\[
h_0i_0 \ =\ \alpha_0\ , \qquad\quad \lambda_X h'_{I'} \ =\ h_0\underbrace{\lambda_B}_{\tilde{i}s}
\tag*{(6.3.16)}
\]
% foot: Version 9 février 2019